% SOURCE_AUTHORITY: sga5_fr_workpass.tex, lines 8421--8433 (LF-normalized unit).
% SOURCE_SHA256: 791F4EFFC5E02832D5D77ED03518C8156D6F07E4C8238B03545DB93D883FBB28
Os lemas 1.2.4.1 e 1.2.4.2 permitem enunciar a proposição seguinte, menos precisa:

\begin{statement}{Proposição 1.2.5}
A categoria dos \(\mathbb Z_\ell\)-feixes constantes torcidos construtíveis sobre um pré-esquema localmente noetheriano conexo \(X\) é equivalente à categoria dos \(\mathbb Z_\ell\)-módulos de tipo finito sobre os quais \(\pi_1(X)\) atua continuamente para a topologia \(\ell\)-ádica.
\end{statement}

\begin{statement}{Observação}
Sejam \(F\) um \(\mathbb Z_\ell\)-feixe constante torcido construtível e \(M\) o \(\mathbb Z_\ell\)-módulo de tipo finito correspondente, dotado de uma ação contínua de \(\pi_1\). Para que \(F\) seja localmente constante, é claro que basta que \(\pi_1\) atue sobre \(M\) através de um quociente finito, o que nem sempre se verifica (1.2.2). Isso é até mesmo necessário quando \(X\) é \emph{normal}. Para vê-lo, após decompor \(X\) em componentes conexas, pode-se reduzir ao caso em que é íntegro. Afirmo então que, denotando por \(K\) seu corpo de funções racionais, se existe uma extensão separável finita \(K'\) de \(K\) tal que a imagem inversa \(F_{K'}\) de \(F\) sobre \(\Spec(K')\) seja constante, então o grupo \(\pi_1(X)\) atua sobre \(M\) através de um grupo finito. Com efeito (SGA 1 V 8.2), \(\pi_1(K)\to \pi_1(X)\) é sobrejetivo, e a hipótese significa que o núcleo do composto
\[
\pi_1(K)\longrightarrow \pi_1(X)\longrightarrow \Aut(M)
\]
é de índice finito. O mesmo acontece, portanto, com o núcleo de \(\pi_1(X)\to\Aut(M)\).
\end{statement}
